% Gemini theme: https://github.com/anishathalye/gemini
% Patched for pdflatex (no fontspec, no Raleway/Lato).
\documentclass[final]{beamer}

% ==== Paper size: 36 in wide; height tuned to fit boxed content ====
% 36 in = 91.44 cm; 37 in = 93.98 cm
\usepackage[size=custom,width=91.44,height=93.98,scale=1.15]{beamerposter}
\usetheme{gemini}
\usecolortheme{gemini}

% ==== Layout ====
\newlength{\sepwidth}
\newlength{\colwidth}
\setlength{\sepwidth}{0.03\paperwidth}
\setlength{\colwidth}{0.46\paperwidth}
\newcommand{\separatorcolumn}{\begin{column}{\sepwidth}\end{column}}

% ==== Section titles: single thick colored rule above for clean separation ====
\setbeamerfont{block title}{family=\sffamily,size=\Large,series=\bfseries}
\newcommand{\blocktitlerule}{%
  {\color{darkblue}\rule{\linewidth}{1.6pt}}\par%
}
\setbeamertemplate{block begin}{%
  \vspace*{0.4ex}%
  \blocktitlerule
  \vspace*{0.6ex}%
  {\usebeamerfont{block title}\color{darkblue}\insertblocktitle\par}%
  \vspace*{0.6ex}%
  \usebeamerfont{block body}%
  \begin{beamercolorbox}[colsep*=0ex]{block body}%
  \justifying%
  \setlength{\parskip}{1ex}%
  \vskip0.3ex%
}
\setbeamertemplate{block end}{\end{beamercolorbox}\vskip0pt\vspace*{1.4ex}}
\setbeamertemplate{block alerted begin}{%
  \vspace*{0.4ex}%
  \blocktitlerule
  \vspace*{0.6ex}%
  {\usebeamerfont{block title}\color{darkblue}\insertblocktitle\par}%
  \vspace*{0.6ex}%
  \usebeamerfont{block body}%
  \begin{beamercolorbox}[colsep*=0ex]{block body alerted}%
  \justifying%
  \begin{adjustwidth}{1ex}{1ex}%
  \setlength{\parskip}{1ex}%
  \vskip0.3ex%
}
\setbeamertemplate{block alerted end}{\end{adjustwidth}\vskip1ex\end{beamercolorbox}\vskip0pt\vspace*{1.6ex}}

% ==== Packages ====
\usepackage{booktabs}
\usepackage{array}
\usepackage{xcolor}
\usepackage{tikz}
\usetikzlibrary{positioning, arrows.meta, shapes.geometric}
\usepackage{graphicx}
\usepackage{listings}

\definecolor{cmured}{RGB}{196,18,48}
\definecolor{accent}{RGB}{30,90,160}
\definecolor{good}{RGB}{20,130,60}
\definecolor{lightgray}{RGB}{240,240,240}

\lstdefinestyle{cppstyle}{
  basicstyle=\ttfamily\small,
  keywordstyle=\color{accent}\bfseries,
  commentstyle=\color{gray}\itshape,
  stringstyle=\color{good},
  language=C++,
  showstringspaces=false,
  breaklines=true,
  frame=none,
  backgroundcolor=\color{lightgray},
  xleftmargin=8pt,
  framexleftmargin=8pt,
  aboveskip=1ex,
  belowskip=1ex,
}

% ==== Title ====
\title{Evolving Hardware Prefetchers with LLM-Driven Code Mutation}
\author{Jerry Hu \quad Helen Wang \quad Hao Kang}
\institute[CMU]{Carnegie Mellon University \quad \texttt{\{zhuoranh,jingzhi3,haok\}@andrew.cmu.edu}}

\begin{document}
\begin{frame}[t,fragile]
\begin{columns}[t]
\separatorcolumn

% =========================================================
% LEFT COLUMN
% =========================================================
\begin{column}{\colwidth}

\begin{block}{Motivation}
Hardware prefetchers predict future cache misses to hide DRAM latency. Modern designs (BOP, IPCP, Berti) expose dozens of \emph{numeric thresholds, scoring functions, and classification rules}; the design space is huge.

\textbf{OpenEvolve} is an LLM-driven evolutionary optimizer. We mark a region of the prefetcher's C++ source with \texttt{EVOLVE-BLOCK} and let \texttt{gpt-5-mini} (CMU LiteLLM gateway) propose mutations. Each candidate is built, simulated in ChampSim, scored, and fed back to a multi-island genetic algorithm.
\end{block}

\begin{block}{Pipeline}
\centering
\begin{tikzpicture}[
  node distance=2.4cm and 2.0cm,
  box/.style={rectangle, draw=cmured, very thick, fill=lightgray, rounded corners=4pt,
              minimum width=5.4cm, minimum height=1.8cm, align=center,
              font=\normalsize\bfseries, inner sep=4pt},
  arrow/.style={-{Stealth[length=7mm,width=5mm]}, line width=2pt, draw=gray!75}
]
\node[box] (seed)  {Seed program\\\footnotesize \texttt{EVOLVE-BLOCK}};
\node[box, right=of seed] (llm)   {LLM mutation\\\footnotesize gpt-5-mini};
\node[box, right=of llm]  (build) {Build\\\footnotesize \texttt{make} champsim};
\node[box, below=of build] (sim)   {ChampSim sim\\\footnotesize 200K + 2M instr};
\node[box, below=of llm]   (score) {Score parse\\\footnotesize IPC, useful, useless};
\node[box, below=of seed]  (ga)    {Multi-island GA\\\footnotesize archive top 10};
\draw[arrow] (seed)  -- (llm);
\draw[arrow] (llm)   -- (build);
\draw[arrow] (build) -- (sim);
\draw[arrow] (sim)   -- (score);
\draw[arrow] (score) -- (ga);
\draw[arrow] (ga)    -- (seed);
\end{tikzpicture}

\medskip
\textbf{Fitness:}
\(\textsf{score} = 0.9 \cdot \tfrac{\text{IPC}}{0.471} + 0.1 \cdot \tfrac{\text{useful}}{\text{useful}+\text{useless}}\)
--- a 9$\times$ asymmetric trade-off.\\[0.4ex]
Accuracy losses are tolerated when they unlock larger speedups.
\end{block}

\begin{block}{The Three Evolved Prefetchers}
All sit at the L2 cache (1024 sets $\times$ 8 ways, 10-cycle latency, 32-entry MSHR).

\heading{BOP --- Best Offset Prefetcher}
Tests candidate offsets $o$: if \mbox{$(\text{block} - o)$} was recently accessed, $o$ scores up. The winning offset triggers prefetches.
Knobs: \texttt{SCORE\_MAX}, \texttt{ROUND\_MAX}, \texttt{PREFETCH\_DEGREE}, \texttt{MSHR\_THRESHOLD}, the 26-element \texttt{OFFSET\_LIST}, and \texttt{score\_update()}.

\heading{IPCP --- IP Class Prefetcher}
Classifies each load PC into \textsc{cs} (constant stride), \textsc{cplx} (signature-based complex), \textsc{gs} (global stream), or \textsc{nl} (next-line); each class has its own prefetch degree.

\heading{Berti --- Timing-Aware Delta Prefetcher}
Per load PC, keeps a history of recent block addresses with issue cycles. Ranks observed deltas by \emph{timeliness} (latency vs.\ avg.\ miss latency); the best delta and its multiples are prefetched.
\end{block}

\begin{block}{Simulated Machine \& Trace}
\centering
\begin{tabular}{ll}
\toprule
\textbf{Component} & \textbf{Configuration} \\
\midrule
Core      & 4\,GHz OoO, 352-entry ROB, bimodal BP \\
L1I / L1D & 64$\times$8 (4c) / 64$\times$12 (5c) \\
L2C       & 1024$\times$8, 10c, evolved prefetcher \\
LLC       & 2048$\times$16, 20c, LRU \\
DRAM      & DDR4-3200, 1ch, $t_{\text{CAS}}{=}t_{\text{RCD}}{=}t_{\text{RP}}{=}24$ \\
\bottomrule
\end{tabular}

\medskip
\begin{flushleft}
\textbf{Trace:} \texttt{602.gcc\_s-734B} (SPEC CPU 2017), 200K warmup + 2M sim instr.\\[0.5ex]
\textbf{Baseline IPC (no prefetcher):} 0.471
\end{flushleft}
\end{block}

\begin{block}{Evolution Configuration}
\centering
\begin{tabular}{l l @{\hspace{2.5em}} l l}
\toprule
\texttt{max\_iterations}      & 20         & \texttt{population\_size}    & 30 / island \\
\texttt{primary\_model}       & gpt-5-mini & \texttt{num\_islands}        & 2           \\
\texttt{temperature}          & 0.7        & \texttt{archive\_size}       & 10          \\
\texttt{max\_tokens}          & 8000       & \texttt{elite\_selection}    & 0.20        \\
\texttt{checkpoint\_interval} & 5          & \texttt{exploitation\_ratio} & 0.70        \\
\bottomrule
\end{tabular}
\end{block}

\end{column}

\separatorcolumn

% =========================================================
% RIGHT COLUMN
% =========================================================
\begin{column}{\colwidth}

\begin{block}{What Did the LLM Actually Change?}

\heading{BOP --- best at iteration 6}
Doubled \texttt{PREFETCH\_DEGREE} (1$\to$2), tightened \texttt{MSHR\_THRESHOLD} (0.7$\to$0.6), shortened \texttt{ROUND\_MAX} (100$\to$64), and duplicated small offsets in \texttt{OFFSET\_LIST}.

\heading{IPCP --- best at iteration 1}
A \emph{structural} change: \texttt{classify\_ip()} was made \textbf{non-static} so it could read \texttt{global\_stream\_dir} and route direction-aligned PCs straight into the \textsc{gs} class.
\begin{lstlisting}[style=cppstyle]
if (global_stream_dir != 0) {
  if ((global_stream_dir > 0 && new_stride > 0) ||
      (global_stream_dir < 0 && new_stride < 0)) {
    if (confidence < CS_CONFIDENCE_THRESHOLD ||
        old_class == GS) return GS;
  }
}
\end{lstlisting}

\heading{Berti --- best at iteration 15}
Thresholds tightened (\texttt{TIMELY} $0.50{\to}0.40$, \texttt{LATE} $1.50{\to}1.12$), weights rebalanced (\texttt{TIMELINESS} $1.0{\to}1.8$, \texttt{ACCURACY} $0.5{\to}1.6$), \texttt{PREFETCH\_DEGREE} doubled. Scoring gained nonlinear penalties:
\begin{lstlisting}[style=cppstyle]
late_penalty = pow(late_ratio, 1.5)
             * (1 + confidence/5.0) * ACCURACY_WEIGHT;
score = coverage + timeliness - late_penalty
      + 0.18*sqrt(n) - 0.03*sqrt(n);
\end{lstlisting}
\end{block}

\begin{block}{Headline Results (gcc\_s)}
\centering
\begin{tabular}{l c c c c}
\toprule
\textbf{Variant} & \textbf{IPC} & \textbf{Speedup} & \textbf{Acc.} & \textbf{Score} \\
\midrule
no prefetcher          & 0.471 & 1.00$\times$ & --     & --   \\
\midrule
BOP (baseline)         & 0.821 & 1.74$\times$ & 98.7\% & 1.66 \\
\textbf{BOP evolved}   & \textbf{1.040} & \textbf{2.21$\times$} & 96.5\% & \textbf{2.08} \\
\midrule
IPCP (baseline)        & 0.811 & 1.72$\times$ & 98.7\% & 1.64 \\
\textbf{IPCP evolved}  & \textbf{1.088} & \textbf{2.31$\times$} & 97.4\% & \textbf{2.18} \\
\midrule
Berti (baseline)       & 0.956 & 2.03$\times$ & 96.8\% & 1.92 \\
\textbf{Berti evolved} & \textbf{1.100} & \textbf{2.34$\times$} & 94.7\% & \textbf{2.20} \\
\bottomrule
\end{tabular}

\medskip
Relative IPC gains over hand-tuned: \textbf{BOP +27\%}, \textbf{IPCP +34\%}, \textbf{Berti +15\%}.
\end{block}

\begin{block}{Convergence: Best Score per Iteration}
\centering
\includegraphics[width=0.98\colwidth]{convergence.pdf}

\textbf{Single-trace saturation.} BOP plateaus by iter 6, IPCP by iter 1, Berti by iter 15.
\end{block}

\end{column}

\separatorcolumn
\end{columns}
\end{frame}
\end{document}
